% =============================================================================
% TEMPLATE DISPENSE UNIVERSITARIE — VARIANTE SCHEMATICA / OUTLINE
% Stile "da appunti": definizioni con i due punti, liste annidate,
% sotto-titoli in grassetto, alberi di classificazione (forest).
% Ricalca lo stile delle dispense di biologia scritte a mano dallo studente.
% =============================================================================
\documentclass[11pt, oneside, openany]{book}

% --- Lingua e codifica ---
\usepackage[italian]{babel}
\usepackage[T1]{fontenc}
\usepackage[utf8]{inputenc}
\usepackage{lmodern}

% --- Margini e layout (ottimizzato per lettura a schermo, non per la stampa) ---
% Margini minimi, pagina singola (niente rilegatura né margini specchiati).
\usepackage[
  margin=12mm, headsep=3mm, footskip=6mm,
  includeheadfoot
]{geometry}
\setlength{\headheight}{15pt}
\raggedbottom

% --- Spaziatura e tipografia ---
\usepackage{setspace}
\setstretch{1.15}
\setlength{\parindent}{0pt}
\setlength{\parskip}{4pt}
\usepackage[protrusion=true, expansion=true, babel=true]{microtype}

% --- Titoli capitolo (deve stare prima di hyperref) ---
\usepackage{titlesec}
\titleformat{\chapter}[block]
  {\normalfont\Huge\bfseries}
  {\thechapter\quad}{0pt}{}
\titlespacing*{\chapter}{0pt}{-60pt}{6pt}
% Intestazioni di sezione: nere, standard (niente colori).

% --- Link e segnalibri PDF ---
\usepackage[hidelinks]{hyperref}
\usepackage{bookmark}

% --- Intestazioni e piè di pagina ---
\usepackage{fancyhdr}
\pagestyle{fancy}
\fancyhf{}
\fancyhead[L]{\leftmark}
\fancyhead[R]{\rightmark}
\fancyfoot[C]{\thepage}
\renewcommand{\headrulewidth}{0.4pt}
\renewcommand{\footrulewidth}{0pt}
\fancypagestyle{plain}{%
  \fancyhf{}%
  \fancyfoot[C]{\thepage}%
  \renewcommand{\headrulewidth}{0pt}%
}
\renewcommand{\chaptermark}[1]{\markboth{Cap.~\thechapter:\ #1}{}}
\renewcommand{\sectionmark}[1]{\markright{\thesection.\ #1}}

% --- Matematica e grafica ---
\usepackage{amsmath}
\usepackage{amssymb}
\usepackage{graphicx}

% --- Tabelle ---
\usepackage{booktabs}
\usepackage{array}
\usepackage{tabularx}
\usepackage{multirow}

% --- Liste annidate (strumento principale dello stile schematico) ---
\usepackage{enumitem}
\setlist[itemize]{noitemsep, topsep=2pt}
\setlist[enumerate]{noitemsep, topsep=2pt}
% Nidificazione di default: enumerate 1. -> (a) -> i. ; itemize • -> – -> *

% --- Due colonne per liste lunghe parallele ---
\usepackage{multicol}

% =============================================================================
% COMANDI DI STILE SCHEMATICO
% =============================================================================
% Sotto-titolo in grassetto, non numerato, su riga propria (es. "Struttura", "DNA").
\newcommand{\rubrica}[1]{\par\addvspace{4pt}\noindent\textbf{#1}\par\nobreak}

% Alberi di classificazione: minimali, solo testo + rami (nessun box).
% Uso: \begin{center}\begin{forest} classalbero [Radice [A][B]] \end{forest}\end{center}
\usepackage{forest}
\forestset{
  classalbero/.style={
    for tree={
      grow'=east, parent anchor=east, child anchor=west,
      font=\footnotesize, anchor=west, align=center, edge={-},
      l sep=8mm, s sep=2mm, inner sep=1pt,
    }
  }
}

% --- Metadati --- (sostituire con i dati del corso)
\title{Dispense di \underline{Valutazione funzionale}\\ \underline{e studio della performance del calciatore}}
\author{}
\date{A.A. 2025/2026}

% =============================================================================
\begin{document}

\frontmatter                 % numerazione romana (i, ii, ...) per titolo e indice

\maketitle
\thispagestyle{empty}
\clearpage

\tableofcontents
\clearpage

\mainmatter                  % la numerazione delle pagine riparte da 1

% =============================================================================
% CAPITOLI — aggiungere \input{capitoli/nome_file} man mano
% =============================================================================

% mazzi: 02
\chapter{Il sistema neuromuscolare}
\label{cap:neuromuscolare}

\textbf{Sistema neuromuscolare}: aspetto principale della prestazione, e quindi il presupposto
per capire \textbf{come} e \textbf{che cosa} valutare.

% ---------------------------------------------------------------------------
\section{I fattori limitanti la prestazione sportiva}

La prestazione sportiva si basa sull'\textbf{individualità del singolo atleta}, condizionata da
tre fattori:
\begin{itemize}
  \item \textbf{genotipo}: corredo genetico di partenza, in cui entrano il grado di
    allenabilità e l'età del soggetto, intesa come \textit{età biologica};
  \item \textbf{allenamento}: tipo di lavoro svolto;
  \item \textbf{stato di salute}, a sua volta condizionato da:
    \begin{itemize}
      \item \textbf{infortuni}: nell'arco della carriera il giocatore ha un alto rischio di
        subirne, rischio maggiore per gli eventi successivi al primo;
      \item \textbf{fatica}: incide in maniera considerevole sul sistema neuromuscolare, ed è
        in grado di alterare la prestazione;
      \item \textbf{dieta}.
    \end{itemize}
\end{itemize}

I tre confluiscono sulla prestazione attraverso le leggi della \textbf{fisiologia}, della
\textbf{psicologia} e della \textbf{biomeccanica}: genotipo + allenamento + salute
$\rightarrow$ atleta $\rightarrow$ prestazione sportiva.

% ---------------------------------------------------------------------------
\section{Il sistema nervoso nel controllo del movimento}

Il sistema nervoso interviene nel movimento con quattro sottosistemi:
\begin{itemize}
  \item \textbf{sistema piramidale}: dall'\textbf{area motoria 4} arriva ai muscoli e invia lo
    stimolo per effettuare il movimento; via del movimento volontario;
  \item \textbf{sistema extrapiramidale}: via del movimento automatico e del controllo
    posturale;
  \item \textbf{sistema propriocettivo}: informa le strutture nervose superiori di tutte le
    variazioni che avvengono nel corpo;
  \item \textbf{sistema esterocettivo}: fornisce informazioni sull'ambiente in cui ci si muove.
\end{itemize}

\rubrica{Influenza sul movimento e sulla postura}
\textbf{Movimento} e \textbf{stabilizzazione posturale} dipendono da:
\begin{itemize}
  \item \textbf{performance muscolare}: velocità, forza, potenza, resistenza, accelerazioni e
    decelerazioni;
  \item \textbf{equilibrio posturale};
  \item \textbf{integrità legamentosa e articolare}.
\end{itemize}

Il circuito di controllo procede così:
\begin{enumerate}
  \item \textbf{recettori cutanei}, \textbf{muscolari} e \textbf{articolari} $\rightarrow$
    midollo spinale afferente;
  \item i \textbf{sistemi motori spinali} generano le risposte riflesse $\rightarrow$ midollo
    spinale efferente $\rightarrow$ attività muscolare riflessa;
  \item le informazioni salgono come \textbf{vista e sensibilità esterocettiva},
    \textbf{sensibilità propriocettiva} e \textbf{informazioni labirintiche};
  \item da qui si dividono in due vie:
    \begin{enumerate}
      \item \textbf{sistemi motori corticali} (coscienti) $\rightarrow$ propriocettività
        cosciente, programmazione motoria, tempo di reazione, sincronizzazione neuromuscolare
        $\rightarrow$ controllo della contrazione muscolare volontaria;
      \item \textbf{sistemi motori sopraspinali} (automatici) $\rightarrow$ percezione
        cinestesica cosciente $\rightarrow$ controllo dell'attività muscolare automatica.
    \end{enumerate}
\end{enumerate}

I parametri cinematici che si misurano --- velocità, spostamento, accelerazione, forza, potenza
muscolare, capacità di resistenza --- \textbf{sottendono quindi un'attività del sistema nervoso
a più livelli}, compresa quella dei meccanocettori e dei propriocettori sparsi in muscoli,
tendini, articolazioni e legamenti, che interagiscono con le afferenze visive e vestibolari.
Dietro ai parametri meccanici e biomeccanici misurati ci sono veri e propri \textbf{schemi
neurali}: il punto è decisivo nello studio dell'infortunio, perché qualsiasi infortunio tende ad
alterare il sistema neurale \textbf{in maniera più stabile} di quanto non faccia con le strutture
biomeccaniche.

% ---------------------------------------------------------------------------
\section{La contrazione muscolare}

Le caratteristiche della contrazione sono \textbf{biofisiche}: \textit{biofisico} è il terreno
comune in cui interagiscono le leggi della \textbf{fisica}, che riguardano l'ambiente e il corpo
umano, e quelle della \textbf{biologia}.

Lo stimolo parte dall'\textbf{area motoria 4}, arriva al \textbf{midollo spinale} e da qui,
attraverso l'innervazione dell'\textbf{$\alpha$-motoneurone}, innesca la sequenza che trasforma
l'\textbf{impulso nervoso in energia meccanica}, trasferendo la forza della contrazione al
tendine e quindi all'osso.

\rubrica{Dall'impulso all'accorciamento}
\begin{enumerate}
  \item l'impulso arriva nella \textbf{zona di innervazione} e si propaga come \textbf{potenziale
    d'azione};
  \item attraverso i \textbf{tubuli T} $\rightarrow$ rilascio di $Ca^{2+}$ dal \textbf{reticolo
    sarcoplasmatico};
  \item il $Ca^{2+}$ libera il \textbf{sito attivo dell'actina};
  \item si forma il \textbf{ponte actomiosinico} (\textit{cross-bridge});
  \item scissione di \textbf{ATP} in \textbf{ADP} $\rightarrow$ scorrimento dell'actina sulla
    miosina;
  \item \textbf{accorciamento del sarcomero}; l'accorciamento di tutti i sarcomeri in serie e in
    parallelo produce l'\textbf{accorciamento muscolare}.
\end{enumerate}

In parallelo agiscono i \textbf{sistemi di controllo} --- fusi neuromuscolari e riflesso spinale
in genere --- che inviano informazioni al midollo spinale per intervenire in modo
\textbf{eccitatorio} o \textbf{inibitorio} secondo il recettore coinvolto, e sono quindi in grado
di modulare il movimento.

% ---------------------------------------------------------------------------
\section{L'unità motoria}

\textbf{Unità motoria} (UM): un motoneurone e tutte le fibre muscolari da esso innervate; è la
parte essenziale della contrazione muscolare.

\begin{itemize}
  \item stimolando il motoneurone, tutte le sue fibre si contraggono \textbf{come una sola
    unità};
  \item il numero di fibre per UM varia da poche unità a parecchie migliaia:
    \begin{itemize}
      \item \textbf{UM piccole} $\rightarrow$ movimenti fini e delicati (muscoli extraoculari);
      \item \textbf{UM ampie} $\rightarrow$ muscoli grandi, come i posturali.
    \end{itemize}
\end{itemize}

\subsection{Tipi di unità motorie}

Due tipi fondamentali:
\begin{itemize}
  \item \textbf{soglia di attivazione bassa}: UM piccole, lente e resistenti alla fatica; nei
    \textbf{muscoli posturali}, che restano contratti a lungo;
  \item \textbf{soglia di attivazione alta}: UM grandi, veloci e facilmente affaticabili; nei
    \textbf{muscoli dinamici}, che permettono la prestazione ad alta intensità.
\end{itemize}

\rubrica{Il ruolo dell'\texorpdfstring{$\alpha$}{alfa}-motoneurone}
Le UM si distinguono \textbf{in funzione delle caratteristiche dell'$\alpha$-motoneurone che le
innerva}: tutte e tre le tipologie di fibre hanno la stessa matrice proteica, ed è il
condizionamento imposto dal singolo motoneurone a differenziarle in rosse, intermedie e veloci.

\rubrica{Le tre tipologie}
\begin{enumerate}
  \item \textbf{S} (\textit{slow}) $\rightarrow$ fibre \textbf{SO} (\textit{slow oxidative}),
    dette \textbf{rosse}:
    \begin{itemize}
      \item motoneurone con frequenza di stimolo molto bassa;
      \item metabolismo aerobico $\rightarrow$ difficilmente affaticabili;
      \item scarsa tensione muscolare;
      \item ricche di enzimi, mitocondri e mioglobina $\rightarrow$ colore rosso; abbondanti nei
        muscoli posturali, sempre attivi.
    \end{itemize}
  \item \textbf{FR} (\textit{fast fatigue-resistant}) $\rightarrow$ fibre \textbf{FOG}
    (\textit{fast oxidative-glycolytic}), dette \textbf{intermedie}:
    \begin{itemize}
      \item motoneurone con frequenza d'impulso superiore;
      \item tensione maggiore delle fibre lente;
      \item \textbf{doppio metabolismo}, aerobico e glicolitico.
    \end{itemize}
  \item \textbf{FF} (\textit{fast fatigable}) $\rightarrow$ fibre \textbf{FG} (\textit{fast
    glycolytic}), dette \textbf{veloci}:
    \begin{itemize}
      \item frequenza di scarica elevata degli $\alpha$-motoneuroni;
      \item elevati gradienti di forza: fibre potentissime;
      \item metabolismo glicolitico anaerobico;
      \item facilmente affaticabili.
    \end{itemize}
\end{enumerate}

\subsection{Il reclutamento delle unità motorie}

\textbf{Principio di Henneman} (o principio della dimensione): utilizzando \textbf{carichi
crescenti} e applicando una tensione leggermente superiore al valore d'inerzia, le UM vengono
reclutate in \textbf{ordine crescente di soglia}, dalle più piccole e lente alle più grandi e
veloci.

\begin{enumerate}
  \item carico basso $\rightarrow$ fibre \textbf{lente} (SO);
  \item aumentando il carico $\rightarrow$ fibre \textbf{intermedie} (FOG);
  \item aumentando ancora $\rightarrow$ fibre \textbf{veloci} (FG);
  \item carico intorno all'\textbf{80\% della forza massima} $\rightarrow$ si reclutano tutte le
    fibre del muscolo;
  \item dall'80\% alla \textbf{forza isometrica massima} la tensione viene ulteriormente
    sviluppata \textbf{non} reclutando nuove fibre, ma aumentando la frequenza degli impulsi.
\end{enumerate}

\rubrica{Reclutamento a carico costante}
Se il carico esterno resta costante --- il \textbf{peso del corpo}, quindi un carico
submassimale --- il reclutamento avviene \textbf{in funzione dell'intensità del movimento}:
\begin{itemize}
  \item movimento lento $\rightarrow$ fibre lente;
  \item aumento di intensità e di velocità di contrazione $\rightarrow$ fibre intermedie;
  \item movimenti esplosivi e \textit{sprint} $\rightarrow$ fibre veloci.
\end{itemize}

Ciò \textbf{non} supera il principio di Henneman: a velocità di contrazione altamente elevate i
tempi sono talmente ridotti che le fibre lente non riescono a intervenire nella contrazione.

\rubrica{Intensità dello stimolo e attivazione delle UM}
Al crescere della percentuale di \textit{pool} reclutato aumenta la percentuale di forza tetanica
sviluppata, con la sequenza S $\rightarrow$ FR $\rightarrow$ F$_{(int)}$ $\rightarrow$ FF. I gesti
si collocano lungo questa scala:
\begin{itemize}
  \item \textbf{stazione eretta} e \textbf{cammino}: reclutamento minimo, entro circa il 25\% del
    \textit{pool};
  \item \textbf{corsa}, a 0,6, 1,2, 1,8 e 3\,m/s: reclutamento crescente, fino a poco più del
    50\%;
  \item \textbf{salto}: reclutamento massimo, oltre l'80\% del \textit{pool}, con forza vicina al
    100\% di quella tetanica.
\end{itemize}

% ---------------------------------------------------------------------------
\section{Variabilità delle risposte muscolari}

Un muscolo scheletrico può essere contratto a \textbf{qualsiasi intensità} e con la \textbf{forza
desiderata}. Contrazioni lente o veloci dell'intero muscolo si ottengono agendo su due
meccanismi:
\begin{enumerate}
  \item \textbf{frequenza dello stimolo} inviata alle fibre muscolari;
  \item \textbf{reclutamento}, cioè numero e dimensioni delle UM coinvolte.
\end{enumerate}

Esempio: sollevare una penna e sollevare un tavolo --- stesso gesto, tensioni muscolari
profondamente diverse.

% ---------------------------------------------------------------------------
\section{Tensione interna e tensione esterna}

\begin{itemize}
  \item \textbf{tensione interna}: in un muscolo scheletrico in contrazione, è la tensione
    generata dalle \textbf{miofibrille} all'interno delle fibre; da sola non produce movimento;
  \item \textbf{elementi elastici in serie} (EES): fibre di \textbf{endomisio},
    \textbf{perimisio}, \textbf{epimisio} e \textbf{tendini}, cui la tensione interna viene
    trasferita;
  \item \textbf{tensione esterna}: la tensione degli EES.
\end{itemize}

\rubrica{Comportamento degli EES}
Gli EES si comportano come \textbf{elastici}: inizialmente si allungano facilmente, ma una volta
allungati diventano rigidi e più adatti a trasferire la tensione esterna alla resistenza
applicata. È l'\textbf{irrigidimento} della struttura elastica allungata a permettere il
trasferimento della tensione muscolare all'osso.

\rubrica{Condizione per il movimento}
Il movimento si produce solo se tensione muscolare prodotta $\neq$ resistenza esterna; se le due
sono uguali, non si ha movimento.

% ---------------------------------------------------------------------------
\section{I tipi di contrazione muscolare}

Due regimi:
\begin{itemize}
  \item \textbf{isometrico} (statico): tensione muscolare $\leq$ carico esterno;
  \item \textbf{non-isometrico} (dinamico): si articola in \textbf{concentrico} ed
    \textbf{eccentrico}.
\end{itemize}

In base alla \textbf{strategia} e alla \textbf{velocità di contrazione}, i movimenti dinamici si
distinguono ulteriormente in:
\begin{itemize}
  \item \textbf{isotonico}: muove un carico in maniera costante; la velocità massima si raggiunge
    \textbf{alla fine} della contrazione, non all'inizio;
  \item \textbf{isocinetico}: alcune macchine consentono di lavorare a \textbf{velocità
    costante}, che però non è un movimento naturale.
\end{itemize}

\begin{center}
\begin{forest} classalbero
[Contrazione muscolare
  [Isometrico\\(statico)]
  [Non-isometrico\\(dinamico)
    [Concentrico [Isotonico][Isocinetico]]
    [Eccentrico [Isotonico][Isocinetico]]
  ]
]
\end{forest}
\end{center}

\subsection{Contrazione isometrica}

Il muscolo sviluppa una \textbf{tensione uguale alla resistenza esterna}, trasmessa al tendine
\textbf{senza produrre movimento esterno}.

\begin{itemize}
  \item il \textbf{ventre muscolare si accorcia} comunque, stirando gli elementi elastici in
    serie rappresentati dal tendine;
  \item tensione muscolare = tensione esterna $\rightarrow$ nessun movimento: i capi articolari
    restano nella posizione di origine.
\end{itemize}

\subsection{Contrazione concentrica}

Il muscolo \textbf{si accorcia} sviluppando una tensione \textbf{maggiore della resistenza
esterna} e produce un movimento che avvicina tra loro i capi articolari; la velocità è sempre
positiva.

\begin{itemize}
  \item il \textbf{picco di velocità} viene \textbf{subito dopo il picco di forza}: quando
    compare, si è già nella fase di distensione;
  \item la \textbf{velocità massima} si raggiunge alla fine, quando la forza torna al livello del
    peso corporeo (\textit{body weight});
  \item segue una piccola \textbf{fase di volo}, dopo la quale la velocità decade secondo
    l'accelerazione di gravità.
\end{itemize}

\subsection{Contrazione eccentrica}

Il muscolo si contrae producendo una tensione \textbf{inferiore alla resistenza esterna}: i capi
articolari \textbf{si allontanano} tra loro.

\begin{itemize}
  \item i livelli di forza prodotti in regime eccentrico sono \textbf{superiori alla forza
    massima} espressa in regime isometrico;
  \item si tratta quindi di \textbf{carichi eccessivi}, oltre la forza isometrica massima:
    esercitazione riservata ad atleti molto esperti, che abbiano già utilizzato il carico e
    possiedano una buona tecnica di sollevamento;
  \item impiego soprattutto nell'\textbf{atletica leggera}, dove picchi elevati di forza e
    potenza sono legati alla prestazione; nei \textbf{giochi di squadra} molto meno consigliabile.
\end{itemize}

% ---------------------------------------------------------------------------
\section{Ciclo stiramento-accorciamento e movimento pliometrico}

\textbf{Ciclo stiramento-accorciamento} (\textit{stretch-shortening cycle}, SSC): sequenza in cui
a una fase di stiramento (\textit{stretch}) segue immediatamente una fase di accorciamento
(\textit{shortening}); corrisponde alla \textbf{contrazione eccentrico-concentrica}.

\begin{itemize}
  \item è una \textbf{strategia evolutiva}, il movimento naturale per eccellenza;
  \item nella \textbf{fase eccentrica} muscolo ed elementi elastici in serie accumulano energia
    elastica $\rightarrow$ restituita nella fase concentrica successiva;
  \item prestazione \textbf{molto più efficiente} di qualsiasi altro tipo di contrazione dinamica.
\end{itemize}

\rubrica{Il tempo di accoppiamento}
\textbf{Tempo di accoppiamento}: tempo che indica il passaggio da una fase all'altra del ciclo.
\begin{itemize}
  \item deve essere il \textbf{più breve possibile};
  \item altrimenti l'energia elastica si disperde in calore e il \textbf{riflesso da stiramento
    non avviene} nella fase concentrica successiva.
\end{itemize}

\rubrica{Movimento pliometrico}
\textbf{Movimento pliometrico}: ciclo di contrazione eccentrico-concentrico caratterizzato
dall'\textbf{impatto con il suolo}.
\begin{itemize}
  \item per essere definito \textbf{regime pliometrico} il \textbf{tempo di contatto al suolo}
    deve realizzarsi in tempi brevissimi;
  \item differenza rispetto al \textit{counter movement jump}, dove il tempo di contatto non è
    rilevante: qui è \textbf{nel tempo di contatto} che si esprime tutta l'azione del sistema
    neuromuscolare, e quindi tutto lo sviluppo di potenza.
\end{itemize}

\rubrica{Modello meccanico}
Ruota e corpo rigido che ruota attorno a uno spigolo; \textbf{modello massa-molla}.

% ---------------------------------------------------------------------------
\section{La stiffness muscolare}

\textbf{Stiffness}: in campo fisiologico indica forza, resistenza, densità e rigidità dei
\textbf{tendini} e del \textbf{tessuto connettivo} del muscolo.

\begin{itemize}
  \item maggiore la \textit{stiffness} di questi tessuti $\rightarrow$ maggiore l'energia
    immagazzinabile durante un movimento eccentrico, poi restituita e liberata nella fase
    concentrica;
  \item è la \textbf{capacità reattiva elastica} che un muscolo produce per eseguire contrazioni
    pliometriche subito dopo il \textbf{prestiramento muscolare};
  \item permette di \textbf{restituire velocemente} l'energia elastica accumulata nella
    deformazione avvenuta durante il breve contatto con il terreno.
\end{itemize}

\rubrica{Stiffness e tempo di contatto}
Più breve il tempo di contatto con il terreno, maggiore il grado di \textit{stiffness} del
soggetto e quindi la sua capacità di esercitare forza nel più breve tempo possibile: un vantaggio
determinante per la prestazione. Interviene in tutte le prestazioni di \textbf{corsa} e di
\textbf{salto} --- salto in lungo, \textit{sprint} --- e soprattutto nella \textbf{fase lanciata},
dove il sistema deve sviluppare alta tensione di forza in tempi di contatto sempre più brevi.

\rubrica{Il ruolo del terreno}
La \textit{stiffness} riguarda \textbf{anche il terreno}: allenarla su terreno soffice non ha
senso. Caso classico l'\textbf{allenamento sulla sabbia}, dove la \textit{stiffness} è assente
$\rightarrow$ degrado non solo della prestazione, ma delle stesse caratteristiche neuromuscolari
del calciatore, che non è più in grado di sviluppare tensioni muscolari nel più breve tempo
possibile.

% ---------------------------------------------------------------------------
\section{L'integrazione sensomotoria}

\rubrica{Recettori muscolari}
\begin{itemize}
  \item \textbf{fusi neuromuscolari} (\textit{muscle spindles}): sensibili alla lunghezza; una
    volta stimolati sono \textbf{eccitatori} nei confronti dell'$\alpha$-motoneurone;
  \item \textbf{organi tendinei del Golgi} (GTO, \textit{Golgi tendon organs}):
    \textbf{inibitori} dell'$\alpha$-motoneurone;
  \item \textbf{corpuscoli del Pacini} (\textit{paciniform corpuscles}): sensibili alle
    variazioni anche di velocità;
  \item \textbf{terminazioni nervose libere} (\textit{free nerve endings}): legate ai livelli
    biochimici locali del muscolo.
\end{itemize}

Nel ciclo stiramento-accorciamento l'aumento di prestazione deriva dall'azione congiunta del
\textbf{riuso dell'energia elastica} e dell'azione \textbf{eccitatoria dei fusi}, mentre i
\textbf{GTO} esercitano il controllo inibitorio.

\rubrica{Il riflesso spinale e il \texorpdfstring{$\gamma$}{gamma}-loop}
\begin{enumerate}
  \item il muscolo \textbf{si allunga} nella fase eccentrica $\rightarrow$ i fusi neuromuscolari
    inviano impulsi che innescano un \textbf{riflesso spinale} (\textit{loop});
  \item il riflesso va a innervare ulteriormente l'\textbf{$\alpha$-motoneurone};
  \item a questa attività si somma sempre l'\textbf{attività discendente} del sistema piramidale,
    cioè volontario;
  \item la prestazione ne risulta aumentata attraverso l'\textbf{azione $\gamma$}, o
    \textbf{$\gamma$-loop}.
\end{enumerate}

L'informazione ha un \textbf{effetto immediato} a livello spinale, ma prosegue nelle vie nervose
superiori --- midollo spinale $\rightarrow$ cervelletto $\rightarrow$ area somatosensoriale
$\rightarrow$ area motoria 4 --- per essere integrata nell'attività volontaria. Vi confluiscono
\textbf{recettori muscolari}, \textbf{recettori cinestesici} e i recettori della sensibilità
cutanea: \textbf{corpuscoli di Meissner} (tatto) e \textbf{terminazioni nervose libere} (dolore e
temperatura).

\rubrica{Fatica nel ciclo stiramento-accorciamento}
La \textit{stiffness} si esprime dentro un sistema complesso, in cui all'attività discendente
dell'area motoria si affiancano i \textbf{sistemi riflessi}. Con la \textbf{fatica} l'azione
riflessa si riduce, e con essa la qualità della \textit{stiffness}:
\begin{enumerate}
  \item deterioramento della funzione muscolare;
  \item ridotta tolleranza all'impatto;
  \item perdita del potenziale di energia elastica;
  \item aumento del lavoro durante la fase di spinta (\textit{push-off}).
\end{enumerate}

\textbf{Vie discendenti} e \textbf{midollo spinale}, tramite le afferenze \textbf{Ia},
\textbf{Ib} e di \textbf{gruppo III e IV}, agiscono sul riflesso e sulla \textit{stiffness} di
anca, ginocchio e caviglia $\rightarrow$ performance; il \textbf{danno muscolare}
(\textit{muscle damage}) interviene su questo circuito.

% mazzi: 03
\chapter{I presupposti scientifici della valutazione funzionale}

\textbf{Presupposti scientifici}: gli aspetti su cui si fonda la valutazione funzionale, e quindi
il punto di partenza per sviluppare le \textbf{metodologie} con cui indagare le caratteristiche
fisiche e motorie ritenute importanti.

% ---------------------------------------------------------------------------
\section{Obiettivi della valutazione funzionale}

\begin{enumerate}
  \item \textbf{indagine sulle qualità fisiche} --- organiche, neuromuscolari, metaboliche --- che
    influenzano la prestazione;
  \item \textbf{controllo ed ottimizzazione dell'allenamento}, in termini di volume e intensità:
    \begin{itemize}
      \item carichi di lavoro eccessivi per un soggetto possono risultare ottimali per un altro
        $\rightarrow$ occorre passare attraverso test diversi per indirizzare il lavoro;
      \item il lavoro va indirizzato in modo \textbf{individuale}, secondo l'esigenza del singolo
        giocatore, e \textbf{specifico}, secondo la disciplina e il ruolo ricoperto;
    \end{itemize}
  \item \textbf{diagnosi funzionale}, cioè controllo degli \textbf{adattamenti biologici}:
    \begin{itemize}
      \item l'allenamento va inteso come \textbf{programmazione di stimoli meccanici};
      \item \textbf{variazione ambientale} $\rightarrow$ una struttura biologica viene sottoposta a
        stimoli meccanici, di cui occorre poi controllare le \textbf{risposte adattive};
    \end{itemize}
  \item \textbf{ricerca ed identificazione del talento}:
    \begin{itemize}
      \item il docente osserva però che sarebbe più corretto parlare di \textbf{inclinazioni
        attitudinali} verso alcuni sport o alcuni ruoli;
      \item il talento non ha bisogno di essere valutato: risalta subito, si fa notare;
      \item i test servono semmai, per tutti coloro che non sono talenti, a individuare le
        inclinazioni verso determinate discipline o determinati ruoli;
    \end{itemize}
  \item \textbf{ricerca scientifica}: unico modo per sviluppare \textbf{protocolli di ricerca} sul
    \textbf{modello di prestazione} e sui \textbf{metodi di allenamento}, per capire se abbiano o
    meno validità.
\end{enumerate}

% ---------------------------------------------------------------------------
\section{Il modello funzionale della prestazione}

Per arrivare al modello funzionale della prestazione esistono due strade.

\begin{enumerate}
  \item \textbf{valutazione dei parametri fisiologici} misurabili durante la prestazione:
    \begin{itemize}
      \item è la via più complicata, benché la tecnologia stia aiutando sempre di più: oggi si
        dispone di strumenti \textbf{miniaturizzati}, tali da non influenzare la prestazione;
      \item resta difficile negli \textbf{sport da contatto} come il calcio, o comunque in partita
        per alcuni parametri;
    \end{itemize}
  \item \textbf{valutazione delle caratteristiche fisiologiche} dell'atleta che pratica quella
    disciplina sportiva:
    \begin{itemize}
      \item si fonda sul fatto che la disciplina praticata condiziona in maniera specifica le
        capacità fisico-organiche del giocatore;
      \item è probabilmente l'unica strada praticabile negli sport da contatto, quindi anche nel
        calcio.
    \end{itemize}
\end{enumerate}

Ne consegue l'\textbf{identificazione del modello di prestazione} attraverso la valutazione delle
caratteristiche \textbf{metaboliche} e \textbf{neuromuscolari} tipiche dell'atleta.

% ---------------------------------------------------------------------------
\section{Classificazione delle discipline sportive}

Attraverso le \textbf{richieste organico-funzionali in gara} si identifica il tipo e la maggior
incidenza di ciascuna di esse nel determinare la prestazione.

\begin{itemize}
  \item discipline a impegno \textbf{prevalentemente aerobico};
  \item discipline a impegno \textbf{aerobico-anaerobico massivo};
  \item discipline a impegno \textbf{prevalentemente anaerobico};
  \item discipline a impegno \textbf{aerobico-anaerobico alternato};
  \item discipline di \textbf{potenza};
  \item discipline di \textbf{destrezza}.
\end{itemize}

\rubrica{Limite di queste classificazioni}
La letteratura scientifica è piena di classificazioni di questo tipo, ma far rientrare il
\textbf{calcio} in una sola di esse è \textbf{improponibile}: a eccezione forse della prima
(impegno prevalentemente aerobico), il calcio si ritrova sostanzialmente in tutte le altre.

% ---------------------------------------------------------------------------
\section{Definizione di test}

\textbf{Test}: unità essenziale della valutazione funzionale motoria.

\begin{itemize}
  \item la richiesta di un \textbf{compito motorio specifico}, in relazione alla sua modalità di
    esecuzione, coinvolgerà capacità fisiche specifiche;
  \item la valutazione si realizza attraverso la misurazione diretta o indiretta di parametri
    fisici, biomeccanici e/o biologici connessi:
    \begin{itemize}
      \item al compito motorio richiesto;
      \item alla sua modalità di esecuzione;
      \item allo strumento di valutazione utilizzato;
    \end{itemize}
  \item le procedure di esecuzione e di misurazione ne definiscono il \textbf{protocollo}.
\end{itemize}

\rubrica{Esempio: la forza esplosiva degli arti inferiori}
Volendo misurare la forza esplosiva degli arti inferiori, l'unità essenziale --- il test ---
potrebbe essere il \textbf{salto verticale}, realizzabile in due modalità differenti:
\begin{itemize}
  \item con le \textbf{mani ai fianchi}, per evitare l'influenza della componente coordinativa;
  \item con le \textbf{braccia libere}.
\end{itemize}

\rubrica{Che cosa deve contenere il protocollo}
Nel protocollo devono rientrare entrambi gli elementi seguenti, perché è la metodica a rendere i
dati confrontabili tra loro:
\begin{enumerate}
  \item la \textbf{modalità di esecuzione} del gesto;
  \item il \textbf{tipo di strumentazione} utilizzata: usare una \textbf{piattaforma di forza} o un
    \textbf{accelerometro} cambia il protocollo, e possono cambiare anche i risultati.
\end{enumerate}

Ciò vale sia per il confronto \textbf{longitudinale} (pre-post del singolo giocatore) sia per
quello \textbf{trasversale} (tra gruppi di giocatori).

\subsection{Caratteristiche dei test}

I test si dividono in funzione del luogo in cui vengono effettuati, in \textbf{test da
laboratorio} e \textbf{test da campo}; ciascuno dei due si articola a sua volta in generali e
specifici, e tutte le categorie confluiscono nella valutazione del \textbf{sistema metabolico} e
del \textbf{sistema neuromuscolare}. Le associazioni prevalenti sono laboratorio $\rightarrow$
generali e campo $\rightarrow$ specifici.

\begin{center}
\begin{forest} classalbero
[Test
  [Da laboratorio [Generali][Specifici]]
  [Da campo [Generali][Specifici]]
]
\end{forest}
\end{center}

\rubrica{Test da laboratorio}
\begin{itemize}
  \item molto in uso negli anni precedenti, quando le strumentazioni erano ingombranti e
    difficilmente trasportabili;
  \item la valutazione era perciò legata alla logistica dello spazio del laboratorio;
  \item perdevano di specificità: davano informazioni importanti dal punto di vista generale, ma
    poco specifiche;
  \item oggi si stanno occupando di situazioni particolari, riguardanti più l'aspetto
    \textbf{clinico} che quello \textbf{prestativo}.
\end{itemize}

\rubrica{Test da campo}
\begin{itemize}
  \item risultano molto più specifici, perché permettono di utilizzare il gesto più vicino a
    quello della prestazione;
  \item lo strumento un tempo disponibile sul campo era solamente il \textbf{cronometro}; oggi si
    dispone di strumentazioni scientifiche valide, capaci di dare informazioni specifiche sia sul
    sistema neuromuscolare sia sul sistema metabolico.
\end{itemize}

% ---------------------------------------------------------------------------
\section{I presupposti scientifici di un test}

Un test deve possedere quattro requisiti:
\begin{enumerate}
  \item \textbf{validità}: capacità di misurare caratteristiche specifiche; un test è valido se
    misura \textbf{quella} grandezza biologica o capacità fisica che ci si è prefissati di
    valutare;
  \item \textbf{riproducibilità}: attendibilità di risultati simili in prove successive;
  \item \textbf{obiettività}: non deve essere influenzato dall'\textbf{operatore} e/o da
    \textbf{fattori ambientali}; se questi influissero, i risultati non riguarderebbero più
    l'atleta ma altri eventi $\rightarrow$ conclusioni errate;
  \item \textbf{specificità}: solo attraverso di essa si può programmare un allenamento adeguato e
    indagare la capacità fisica che davvero si avvicina al \textbf{modello di prestazione}.
\end{enumerate}

% ---------------------------------------------------------------------------
\section{La validità}

\subsection{Validità sostanziale o di costrutto}

\textbf{Validità sostanziale} (o \textbf{di costrutto}): validità di primaria importanza, che
identifica l'abilità di un test a rappresentare i presupposti teorici che spiegano alcuni aspetti
derivanti dalla conoscenza e dalle osservazioni.

\begin{itemize}
  \item rappresenta il \textbf{limite della misura} di un test nella valutazione delle
    caratteristiche fisiche per le quali è stato progettato;
  \item p.\,es.\ per misurare la forza esplosiva degli arti inferiori occorre un test in grado di
    restituire \textbf{quel} risultato in maniera specifica, senza compromessi.
\end{itemize}

\subsection{Validità formale o apparente}

\textbf{Validità formale} (o \textbf{apparente}): validità di secondaria importanza, che identifica
il \textbf{grado di apparenza esterna} che un test possiede nel misurare le caratteristiche fisiche
e biologiche che si propone di misurare.

\begin{itemize}
  \item è un presupposto di tipo \textbf{qualitativo}: una caratteristica informale e non
    quantitativa;
  \item il soggetto che esegue il test lo riconosce come vicino alla prestazione $\rightarrow$ più
    motivato a eseguire il test o la batteria di test;
  \item la \textbf{motivazione} è essa stessa un aspetto fondamentale della validità del dato: deve
    essere altissima, altrimenti il dato non ha alcun valore;
  \item p.\,es.\ misurare un parametro metabolico di un calciatore usando una \textit{bike} ha una
    validità apparente bassissima: il giocatore capisce da sé che quel tipo di valutazione non ha
    nesso con la sua prestazione $\rightarrow$ la motivazione ne risente.
\end{itemize}

\subsection{Validità di contenuto}

\textbf{Validità di contenuto}: identifica l'abilità di un test e/o di una batteria di test di
misurare tutte le abilità motorie necessarie per una determinata disciplina sportiva.

\begin{itemize}
  \item un \textbf{solo test} non basta a valutare tutte le capacità motorie necessarie allo
    sviluppo della prestazione $\rightarrow$ occorre una \textbf{batteria} di test, che dia un
    quadro più completo di come il giocatore sviluppa la propria prestazione (forza esplosiva,
    \textit{stiffness}, capacità anaerobiche, \dots);
  \item l'importanza data alle abilità motorie indagate deve essere correlata alle richieste del
    compito della disciplina sportiva praticata.
\end{itemize}

\subsection{Validità del criterio di riferimento}

\textbf{Validità del criterio di riferimento}: mette in correlazione la misura di un test con
altri test che misurano la stessa abilità. Se ne distinguono quattro tipi.

\begin{enumerate}
  \item \textbf{validità concorrente}: grado della misura di un test e della strumentazione
    utilizzata, associato (stimato statisticamente) a quella di un altro test e della relativa
    strumentazione, già accettati (ritenuti validi) per la misura della stessa abilità;
    \begin{itemize}
      \item si ha quando due test si avvicinano l'uno all'altro sia per la \textbf{capacità motoria
        indagata} sia per la \textbf{strumentazione utilizzata};
    \end{itemize}
  \item \textbf{validità convergente}: \textbf{alta correlazione} tra le misure di un test e quelle
    di un test riconosciuto come \textit{gold standard};
    \begin{itemize}
      \item p.\,es.\ il \textbf{salto in lungo}, considerato espressione della forza esplosiva, ha
        validità convergente con il \textbf{salto verticale};
    \end{itemize}
  \item \textbf{validità predittiva}: limite della misura di un test in relazione alla misura
    ottenuta nello stesso da \textbf{atleti di alto livello} praticanti la stessa disciplina
    sportiva, cioè da coloro che rappresentano il \textbf{modello di prestazione} $\rightarrow$ la
    misura può essere predittiva della prestazione futura;
  \item \textbf{validità discriminante}: abilità di un test di distinguersi tra due differenti
    costrutti;
    \begin{itemize}
      \item si evidenzia con una scarsa correlazione tra i risultati del test rispetto a quelli
        di un altro costrutto;
      \item buona validità discriminante in una batteria $\rightarrow$ evita inutili spese di
        tempo, energie e risorse nella somministrazione di test altamente correlati tra loro;
      \item i test di una batteria devono misurare in modo indipendente le abilità motorie
        componenti la prestazione;
      \item p.\,es.\ \textbf{salto in lungo} e \textbf{salto in alto} non hanno motivo di coesistere
        nella stessa batteria: indagano la medesima capacità $\rightarrow$ sovrapposizione inutile
        di dati, tempo ed energie.
    \end{itemize}
\end{enumerate}

% ---------------------------------------------------------------------------
\section{La riproducibilità}

\textbf{Riproducibilità}: misura del \textbf{grado di consistenza} o di riproducibilità di un test.

\begin{itemize}
  \item se un atleta possiede un'abilità indagata dal test, i risultati non devono cambiare in
    due misure successive;
  \item il test è riproducibile solo se si ottengono le stesse misure con \textbf{scarsissima
    variabilità} in due momenti successivi di somministrazione;
  \item una variabilità fra misure eseguite a distanza di un giorno che non sia imputabile ad
    alcun fattore esterno, ma solo all'esecuzione del test, rende il test non riproducibile;
  \item un test deve essere \textbf{valido e riproducibile}, perché un'elevata variabilità non
    avrebbe senso.
\end{itemize}

\rubrica{Riproducibilità test-retest}
\textbf{Test-retest}: amministrare il test due volte allo stesso gruppo di atleti ed effettuare
la correlazione statistica tra le due misure.

\subsection{Validità e riproducibilità a confronto}

I due requisiti sono \textbf{indipendenti}, e se ne danno quattro combinazioni.

\begin{enumerate}
  \item \textbf{valido ma non riproducibile}: la \textbf{media} delle misure ripetute si avvicina
    al valore reale, ma le singole misure presentano una variabilità molto alta;
  \item \textbf{riproducibile ma non valido}: le misure ripetute sono tutte vicine o uguali tra
    loro --- scarsissima variabilità --- ma lontane dal valore reale atteso;
  \item \textbf{né valido né riproducibile}: le misure sono distanti tra loro e distanti dal valore
    di riferimento;
  \item \textbf{valido e riproducibile}: non solo c'è scarsa variabilità tra le misure prese
    singolarmente, ma tutte si avvicinano al valore reale dell'abilità indagata.
\end{enumerate}

% ---------------------------------------------------------------------------
\section{La specificità}

\textbf{Specificità di una valutazione}: caratteristica di un test di basarsi, \textbf{laddove
possibile}, su condizioni simili a quelle ambientali e fisiche della competizione, cioè sulle
dinamiche dell'azione e del gesto sportivo.

\begin{itemize}
  \item il test è \textbf{specifico} se riproduce e soddisfa le condizioni della prestazione in
    termini di \textbf{coerenza biomeccanica}, coordinativa e metabolica.
\end{itemize}

\rubrica{Test funzionali specifici}
I \textbf{test funzionali specifici} mirano alla diagnosi di qualità fisiologiche, neuromuscolari
e biomeccaniche tipiche di una disciplina sportiva.

\rubrica{Esempio: le catene cinetiche nel gesto del calciatore}
Trattandosi di movimenti alquanto complessi, saper organizzare test specifici permette di
distinguere le azioni diverse dei due arti nello stesso gesto:
\begin{itemize}
  \item \textbf{arto calciante} $\rightarrow$ catena cinetica \textbf{aperta};
  \item \textbf{arto di appoggio} $\rightarrow$ catena cinetica \textbf{chiusa}, con la funzione
    completamente diversa di mantenere l'equilibrio del corpo.
\end{itemize}

La distinzione risulta interessante non tanto ai fini della prestazione, quanto nella \textbf{fase
di recupero post-infortunio}.

% mazzi: 04
\chapter{Validità degli strumenti di misura}

Dopo i presupposti scientifici della valutazione funzionale, un altro aspetto fondamentale è la
\textbf{validità degli strumenti di misura} impiegati nella valutazione. La vera difficoltà non sta
tanto nel somministrare il test, quanto nel \textbf{ridurre il più possibile gli errori} --- di
metodo, di esecuzione e di misurazione --- che in alcune situazioni sono quasi fisiologici.

% ---------------------------------------------------------------------------
\section{Accuratezza e precisione di una misura}

\subsection{Accuratezza}

\textbf{Accuratezza}: quantifica la \textbf{distorsione} in una misura causata da errori di
misurazione.

\begin{itemize}
  \item è quantificata come la \textbf{differenza tra un valore vero e un valore atteso /
    osservato};
  \item una misurazione effettuata con alta precisione avrà sempre un \textbf{piccolo errore}, cioè
    una piccola distorsione dal valore vero;
  \item viene normalmente indicata come l'\textbf{errore massimo} che può esistere in una
    misurazione, ma in realtà ciò che si quantifica è l'\textbf{imprecisione}, misurando la
    \textbf{deviazione tra le misure};
  \item non riguarda quindi la \textbf{media} tra le misure, bensì la \textbf{deviazione standard}
    delle misure, cioè la \textbf{variabilità} tra una misura e l'altra.
\end{itemize}

\subsection{Precisione}

\textbf{Precisione}: è la \textbf{differenza tra un valore osservato e un valore medio atteso}.

\begin{itemize}
  \item se un sistema misura qualcosa che dovrebbe restare invariata, una \textbf{variazione nei
    valori misurati} indica una \textbf{mancanza di precisione};
  \item in genere viene quantificata utilizzando la \textbf{deviazione standard} delle misure
    \textbf{ripetute} della stessa quantità;
  \item precisione e accuratezza sono a volte confuse, ma sono \textbf{quantitativi distinti}.
\end{itemize}

\rubrica{Esempio: il podobarografo a sensore singolo}
Misurando il \textbf{picco di pressione sotto i piedi} durante il passo con un podobarografo
costituito da un \textbf{solo sensore}:
\begin{itemize}
  \item il sistema sarà \textbf{molto preciso}, ma solo \textbf{in quel punto};
  \item sarà invece \textbf{lontano dalla realtà}, perché si perde tutta l'informazione sulla
    \textbf{distribuzione} della pressione sotto il piede nel suo complesso.
\end{itemize}

\subsection{Combinazioni di accuratezza e precisione}

La figura delle slide illustra le tre permutazioni possibili in funzione del \textbf{bersaglio},
che rappresenta il \textbf{valore vero}.

\begin{enumerate}
  \item[(a)] \textbf{alta precisione e alta accuratezza}: la misurazione centra l'obiettivo in
    maniera completa, con entrambe le caratteristiche;
  \item[(b)] \textbf{alta precisione e bassa accuratezza}: la variabilità delle misurazioni è
    piccolissima, ma l'insieme delle misure è \textbf{lontano dal valore vero};
  \item[(c)] \textbf{bassa precisione e scarsa accuratezza}: non si verifica nessuna delle due
    condizioni --- la variabilità tra le misure è enorme e l'accuratezza è scarsa.
\end{enumerate}

\rubrica{Nota sulla slide}
Nella slide la voce (c) riporta ``bassa precisione e scarsa precisione'': il docente segnala
l'\textbf{errore} e corregge, precisando che il secondo termine sta per \textbf{accuratezza}.

% ---------------------------------------------------------------------------
\section{Grandezze misurabili}

Tutte concorrono, a vario titolo, alla valutazione funzionale delle capacità motorie; le prime due
famiglie derivano più o meno direttamente dalla \textbf{fisica}.

\begin{itemize}
  \item \textbf{grandezze cinematiche}: tempo, distanza, velocità, accelerazione;
  \item \textbf{grandezze dinamiche}: forza, coppia, potenza, lavoro, impulso, quantità di moto;
  \item \textbf{grandezze fisiologiche}: FC, $VO_2max$, lattato ematico, EMGs, temperatura.
\end{itemize}

% ---------------------------------------------------------------------------
\section{Il processo di misurazione strumentale}

Gli strumenti sfruttano la \textbf{variazione di tensione} che si sviluppa in seguito alla
prestazione misurata:

\begin{center}
grandezza misurata $\rightarrow$ trasduttore $\rightarrow$ tensione
\end{center}

\rubrica{Trasduttori}
Potenziometri, elettrogoniometri, encoder, celle di carico, pedana dinamometrica, accelerometri.

\begin{itemize}
  \item p.\,es.\ le \textbf{celle di carico}, sottoposte a forza, subiscono una \textbf{piccola
    deformazione} che crea una \textbf{variazione di tensione}; attraverso il trasduttore questa
    fornisce la grandezza misurata, secondo il tipo di scala e le unità di misura utilizzate;
  \item i valori misurati dalla strumentazione \textbf{contengono normalmente errori}, anche nelle
    strumentazioni più precise.
\end{itemize}

% ---------------------------------------------------------------------------
\section{Proprietà dei trasduttori}

Le seguenti proprietà dei trasduttori possono \textbf{influenzare la validità della misura}.

\subsection{Sensibilità}

\textbf{Sensibilità}: riguarda la \textbf{relazione tra il segnale di uscita e quello di entrata};
è data dal \textbf{rapporto tra la variazione della grandezza in ingresso e la variazione della
grandezza in uscita}.

\begin{itemize}
  \item un trasduttore è \textbf{molto sensibile} quando a una \textbf{piccola variazione
    dell'ingresso} corrisponde una \textbf{grande variazione dell'uscita};
  \item è \textbf{poco sensibile} quando a un'\textbf{elevata variazione del segnale di ingresso}
    l'uscita risponde con una \textbf{piccola variazione};
  \item p.\,es.\ una piattaforma di forza con sensibilità di \textbf{10\,N/v}: ogni 10\,N di forza
    generano una tensione di 1\,V (quindi 20\,N $=$ 2\,V). In questo caso la sensibilità è
    \textbf{ridotta}.
\end{itemize}

\subsection{Errore di non linearità}

\textbf{Errore di non linearità}: teoricamente, quando la grandezza misurata aumenta, anche
l'output di tensione deve \textbf{crescere linearmente}; lo \textbf{scostamento dalla linearità}
rappresenta l'errore dello strumento.

\begin{itemize}
  \item p.\,es.\ la \textbf{cella di carico} sottoposta a una forza esterna si deforma in modo
    \textbf{proporzionale} --- quindi lineare --- all'aumentare della forza applicata; è questa
    relazione a permettere la \textbf{trasformazione dei volt} (la variazione di tensione subita
    dallo \textit{strain gauge}) nella forza esterna applicata;
  \item l'errore è di solito fornito \textbf{direttamente dal costruttore} ed è espresso come
    \textbf{spostamento massimo dalla retta ideale};
  \item in condizioni di linearità lo strumento \textbf{non presenta questo errore}.
\end{itemize}

\subsection{Isteresi}

\textbf{Isteresi}: si verifica quando, per \textbf{valori crescenti} della grandezza in ingresso, la
caratteristica assume un determinato andamento, mentre per \textbf{valori decrescenti} della
medesima grandezza ne disegna uno \textbf{differente}.

\begin{itemize}
  \item un buon trasduttore è caratterizzato da un'isteresi \textbf{praticamente nulla}: la
    variazione all'\textbf{aumentare} della grandezza fisica misurata deve essere \textbf{uguale} a
    quella al suo \textbf{diminuire};
  \item molte strumentazioni presentano quasi sempre una \textbf{leggera isteresi}.
\end{itemize}

\subsection{Cross-talk}

\textbf{Cross-talk}: si verifica solitamente nei trasduttori \textbf{triassiali} (accelerometri,
pedane di forza) quando, esercitando una forza su un \textbf{solo asse} (p.\,es.\ l'asse $z$), si
osserva una \textbf{variazione di segnale anche nelle altre direzioni}.

\begin{itemize}
  \item deve essere \textbf{ridotto al minimo}: altrimenti si crea \textbf{interferenza} e quindi
    una \textbf{variabilità enorme} della strumentazione.
\end{itemize}

\subsection{Risoluzione}

\textbf{Risoluzione}: per ogni strumento di misura vi è un \textbf{limite alla grandezza misurata}
che produce un cambiamento nell'uscita dello strumento; la risoluzione riflette perciò la
\textbf{``finezza''} con cui una misurazione può essere effettuata.

\begin{itemize}
  \item spesso è espressa come \textbf{valore} o come \textbf{percentuale del range di misura
    completo} dello strumento;
  \item con strumenti elettronici, riportare misure con \textbf{molti decimali non indica
    necessariamente un'alta risoluzione}: una grande componente di quei valori può riflettere il
    \textbf{rumore del sistema}.
\end{itemize}

\subsection{Temperatura}

\textbf{Temperatura}: anche la temperatura dell'\textbf{ambiente} in cui si eseguono le misurazioni
può determinare \textbf{variazioni del segnale elettrico in uscita}.

\begin{itemize}
  \item per confrontare misurazioni effettuate in \textbf{tempi diversi} è necessario che siano
    eseguite nelle \textbf{stesse condizioni ambientali};
  \item il motivo è che i \textbf{metalli} che costituiscono gran parte della strumentazione
    elettronica sono sensibili alla temperatura e \textbf{variano le proprie condizioni fisiche} al
    variare di essa;
  \item nella pratica un test svolto sul campo d'estate o a inizio primavera è completamente diverso
    da uno svolto in inverno o a inizio preparazione fisica: di questo aspetto va tenuto conto.
\end{itemize}

% ---------------------------------------------------------------------------
\section{La calibrazione}

\textbf{Calibrazione}: procedura che \textbf{non è diretta a eliminare gli errori}, ma dovrebbe
aiutare a \textbf{tenerli al minimo}; sfrutta un \textbf{modello} che relaziona un \textbf{ingresso
noto} allo strumento di misurazione con un \textbf{output desiderato}.

\begin{itemize}
  \item p.\,es.\ la calibrazione di un \textbf{encoder lineare} consiste nel fornire la misura di
    una \textbf{distanza certa e assoluta}, che la strumentazione metterà in relazione al proprio
    sistema di misura;
  \item questi modelli \textbf{non sono rappresentazioni perfette} del sistema: anche durante
    l'acquisizione dei dati gli errori possono provenire da \textbf{fonti diverse};
  \item gli errori nascono perché le \textbf{condizioni cambiano} rispetto a quando la calibrazione
    è stata eseguita --- p.\,es.\ la \textbf{temperatura dei sensori}; per alcune strumentazioni la
    calibrazione va perciò rifatta \textbf{sistematicamente il giorno stesso del test};
  \item per determinare molti parametri biomeccanici è necessario \textbf{combinare parametri e
    variabili di fonti diverse}: in questo caso gli errori sorgono attraverso la \textbf{propagazione
    degli errori} nella loro combinazione (p.\,es.\ per calcolare la \textbf{quantità di moto} devono
    essere note massa e velocità dell'oggetto);
  \item la buona pratica sperimentale è \textbf{minimizzare gli errori alla fonte}.
\end{itemize}

% ---------------------------------------------------------------------------
\section{Evoluzione e applicazione pratica dell'encoder}

\textbf{Encoder}: \textbf{sensore digitale di posizione} di tipo \textbf{rotativo}.

\begin{itemize}
  \item fornisce una rappresentazione \textbf{temporale}, \textbf{numerica} e \textbf{diretta} di
    \textbf{rotazioni angolari}, cioè la variazione della distanza nel tempo;
  \item attraverso vari meccanismi fornisce anche una \textbf{misura diretta di spostamenti
    rettilinei}.
\end{itemize}

\subsection{Struttura e funzionamento}

È formato da un \textbf{disco} (metallo, ceramica, ecc.) che ruota tramite un'\textbf{asta}
(albero), da un \textbf{LED} (laser) e da un \textbf{rilevatore}.

\begin{itemize}
  \item si tratta di strumentazione \textbf{miniaturizzata}: il disco è molto piccolo e presenta
    lungo la circonferenza dei \textbf{fori} a distanza talmente ridotta tra loro da poterli
    considerare \textbf{lineari};
  \item l'\textbf{emettitore ottico} colpisce il rilevatore in maniera \textbf{alternata}, al
    passaggio dei punti vuoti e di quelli pieni;
  \item può lavorare in due modalità:
    \begin{itemize}
      \item \textbf{in trasparenza}: i segmenti del disco sono alternativamente \textbf{opachi e
        trasparenti};
      \item \textbf{in riflessione}: i segmenti del disco sono alternativamente \textbf{opachi e
        riflettenti};
    \end{itemize}
  \item il \textbf{conteggio del numero di segnali} rilevati dà il numero di passaggi effettuati dal
    disco: \textbf{calibrando} questi passaggi con una \textbf{misura certa} (un metro) si stabilisce
    la \textbf{distanza percorsa} da qualsiasi oggetto a cui l'encoder sia attaccato.
\end{itemize}

\subsection{Applicazione: encoder su macchina isotonica}

Applicando l'encoder al \textbf{pacco pesi} di una \textit{leg extension}, dai tracciati di
\textbf{posizione} e \textbf{velocità} in funzione del tempo si legge l'intero gesto:

\begin{itemize}
  \item \textbf{inizio del movimento}: posizione registrata $=$ 0 e velocità $=$ 0;
  \item \textbf{metà del movimento} (in questo caso $45^\circ$): lo spostamento registrato
    dall'encoder, essendo lineare, corrisponde al \textbf{picco di velocità};
  \item \textbf{massima escursione articolare} ($180^\circ$, gamba distesa): la velocità torna a
    \textbf{zero};
  \item \textbf{ritorno alla fase iniziale}: posizione $=$ 0 e velocità $=$ 0.
\end{itemize}

\rubrica{Dalla posizione ai parametri dinamici}
Avendo la posizione in funzione del tempo si ricava, in cascata:
\begin{itemize}
  \item la \textbf{velocità} (prima misura calcolata);
  \item dalla velocità in funzione del tempo, l'\textbf{accelerazione};
  \item con una \textbf{massa nota}, la \textbf{forza} ($F = m \cdot a$);
  \item da forza e velocità, la \textbf{potenza} ($P = F \cdot v$).
\end{itemize}

Con un sistema semplice si ottengono così, per qualsiasi movimento effettuato su una macchina
isotonica o su un castello, il valore di posizione e \textbf{tutti i parametri cinematici e
dinamici}.

% ---------------------------------------------------------------------------
\section{Grandezza misurata e grandezze calcolate}

Per ridurre l'errore su qualsiasi strumento di misura occorre \textbf{conoscere qual è la grandezza
fisica realmente misurata}: solo così si capisce quali \textbf{calcoli} effettuerà poi il sistema o
il software. Si distinguono perciò una \textbf{grandezza fisica misurata} e le \textbf{grandezze
fisiche calcolate}.

\begin{itemize}
  \item \textbf{encoder}: la grandezza misurata è lo \textbf{spostamento}; si procede per
    \textbf{derivazione} --- $spostamento \rightarrow velocità \rightarrow accelerazione$ (derivata
    prima e derivata seconda);
  \item \textbf{accelerometro}: la grandezza misurata è l'\textbf{accelerazione}; si procede in
    senso inverso, per \textbf{integrazione} --- $accelerazione \rightarrow velocità \rightarrow
    spostamento$ (integrazione prima e seconda).
\end{itemize}

\rubrica{Propagazione degli errori}
\begin{itemize}
  \item l'errore sulla grandezza misurata \textbf{si propaga} alle grandezze calcolate, in funzione
    della \textbf{formula}, del \textbf{livello dell'equazione} e del \textbf{grado di potenza
    dell'algoritmo} utilizzato nelle operazioni successive;
  \item nell'encoder un errore minimo sullo spostamento si propaga sulla velocità e, con
    un'equazione di \textbf{secondo grado}, sull'accelerazione;
  \item questo può rendere le \textbf{misure calcolate scarsamente valide};
  \item da qui l'obiettivo fondamentale di tutte le procedure di valutazione: \textbf{ridurre il più
    possibile l'errore della misurazione} alla fonte.
\end{itemize}

\chapter{Dinamometria isoinerziale}

Con la dinamometria isoinerziale si entra nella \textbf{parte specifica della valutazione},
dopo aver trattato i presupposti scientifici che sottendono questo tipo di metodica.

% ---------------------------------------------------------------------------
\section{Le piattaforme di forza}

\textbf{Piattaforma di forza}: strumentazione tra le più importanti e largamente usate, in uso già
dall'inizio degli anni Settanta.

\begin{itemize}
  \item funziona grazie agli \textit{strain gauge} (estensimetri), normalmente \textbf{quattro},
    posizionati agli angoli della piattaforma;
  \item registra le \textbf{variazioni della forza di reazione del terreno in funzione del tempo};
  \item fornisce informazioni \textbf{vettoriali}: le componenti della forza sui tre assi e i
    momenti rispetto al centro di pressione.
\end{itemize}

\rubrica{Campi di applicazione}
\begin{enumerate}
  \item \textbf{locomozione}: variazioni della forza di reazione durante cammino e corsa, in tutte
    le modalità $\rightarrow$ informazioni sui vettori che attraversano tutto il corpo;
  \item \textbf{controllo posturale}: variazione del centro di massa (o centro di gravità) in
    posizione eretta.
\end{enumerate}

\textbf{Sway path}: tracciato delle variazioni nello spazio e nel tempo del centro di massa durante
il mantenimento della stazione eretta; detto anche ``gomitolo'' per la forma del grafico.
\begin{itemize}
  \item è strettamente connesso alla \textbf{capacità di equilibrio} del soggetto;
  \item il grafico fornisce informazioni non solo sul centro del gomitolo --- relativo al centro di
    massa --- ma anche sui tracciati dell'\textbf{arto inferiore destro e sinistro}.
\end{itemize}

% ---------------------------------------------------------------------------
\section{Applicazioni al controllo posturale}

\subsection{Stimolazione plantare ed equilibrio}

Studio condotto con piattaforma di forza per capire come la \textbf{stimolazione dei meccanocettori
plantari} modifichi il controllo posturale (Annino et al., \textit{Somat \& Motor Res}, 2015).

\rubrica{Strumentazione}
Pedana posturale $\rightarrow$ amplificatore $\rightarrow$ personal computer per l'elaborazione dei
tracciati.

\rubrica{Superfici confrontate}
\begin{enumerate}
  \item \textbf{firm}: superficie rigida e liscia;
  \item \textbf{foam}: superficie morbida, come quella delle solette delle scarpe da \textit{running};
  \item \textbf{FTD} (\textit{firm textured disc}): superficie rigida con protuberanze molto piccole,
    in grado di stimolare i meccanocettori plantari; disco di $\phi$ 17\,mm.
\end{enumerate}

\rubrica{Parametri misurati}
\begin{itemize}
  \item \textbf{CoP}: sway path, cioè la distanza coperta dal baricentro durante la durata del test;
  \item \textbf{$V_{A/P}$}: velocità di oscillazione antero-posteriore;
  \item \textbf{$V_{M/L}$}: velocità di oscillazione medio-laterale.
\end{itemize}
Tutte le misure sono state effettuate sia a \textbf{occhi aperti} sia a \textbf{occhi chiusi}.

\rubrica{Risultati}
\begin{itemize}
  \item la superficie stimolante (FTD) ha \textbf{ridotto lo sway path} sia a occhi aperti sia,
    in modo marcato, a occhi chiusi $\rightarrow$ aumento del controllo posturale;
  \item la riduzione ha riguardato anche le \textbf{velocità di spostamento} in entrambe le
    direzioni, antero-posteriore e laterale, in tutte le condizioni;
  \item la \textbf{superficie soffice} è quella che più destabilizza il controllo posturale: non
    stimolando i meccanocettori plantari riduce le afferenze, e il controllo diventa più complicato.
\end{itemize}

\rubrica{Implicazione allenante}
Le esercitazioni di equilibrio avrebbero probabilmente maggiore efficacia se effettuate su
\textbf{superfici rigide e stimolanti} anziché su superfici soffici --- il contrario di come
l'equilibrio viene comunemente allenato, anche nel calcio.

\subsection{Solette testurizzate e fatica muscolare}

Studio sull'effetto delle \textbf{solette testurizzate} sul controllo posturale in posizione eretta
statica dopo affaticamento degli arti inferiori (Palazzo et al., \textit{J Sports Med Phys Fitness},
2017). Le solette contengono piccoli ciottoli non deformabili di varie dimensioni, capaci di
stimolare un numero maggiore di meccanocettori, e vengono inserite nella scarpa.

\rubrica{Metodo}
\begin{itemize}
  \item \textbf{campione}: 10 giovani adulti sani, età media 26,1 $\pm$ 3,07 anni;
  \item \textbf{condizioni sensoriali}: quattro --- con e senza vista, con e senza stimolazione
    plantare naturale;
  \item \textbf{protocollo}: per ogni condizione un singolo trial di 30\,s, in \textit{pre-fatigue}
    e in \textit{post-fatigue};
  \item \textbf{induzione della fatica}: 60\,s di salti continui;
  \item \textbf{misure}: spostamento del centro di pressione (\textit{CoPDISP}), \textit{sway
    velocity}, velocità di oscillazione antero-posteriore e medio-laterale.
\end{itemize}

\rubrica{Risultati e conclusione}
\begin{itemize}
  \item effetto \textbf{stabilizzante} rispetto alle solette di controllo;
  \item riduzione significativa del \textit{CoPDISP} nella condizione a occhi chiusi pre-fatica;
  \item in post-fatica \textbf{tutti} i parametri posturali sono migliorati, sia con sia senza vista;
  \item conclusione: le solette testurizzate a stimolazione rigida migliorano il \textit{CoPDISP}
    \textbf{indipendentemente dalla vista}, fornendo un'informazione sensoriale completa che
    migliora l'equilibrio in condizioni di fatica.
\end{itemize}

\rubrica{Interpretazione}
Buone afferenze propriocettive possono \textbf{ridurre gli effetti negativi della fatica}, tra cui
i livelli di co-contrazione: si rende più efficiente l'azione del sistema neuromuscolare,
soprattutto nel rapporto tra agonisti e antagonisti.

% ---------------------------------------------------------------------------
\section{Analisi del passo su treadmill}

Per l'analisi del passo della corsa si utilizza un \textbf{treadmill multifunzione} dotato di
piattaforma di pressione (\textit{Zebris FDM-TLR}).

\begin{itemize}
  \item \textbf{sensori}: 5376 sensori capacitivi di nuova generazione;
  \item \textbf{frequenza di campionamento}: 100\,Hz;
  \item \textbf{velocità operativa}: 0--14\,km/h.
\end{itemize}

\rubrica{Grandezze analizzate}
\begin{itemize}
  \item \textbf{forza massima} (N) nelle varie fasi del cammino;
  \item \textbf{pressione} (N/cm$^2$): forza esercitata sull'unità di superficie;
  \item \textbf{lunghezza della falcata} (cm);
  \item \textbf{cadenza dei passi} (step/min);
  \item \textbf{doppia fase statica} (\%);
  \item \textbf{swing phase} (\%): percentuale della fase dell'arto oscillante.
\end{itemize}

\rubrica{Ciclo del passo}
$doppio\ appoggio\ iniziale \rightarrow appoggio\ singolo \rightarrow doppio\ appoggio\ terminale
\rightarrow oscillazione$.
\begin{itemize}
  \item doppia fase statica e fase di oscillazione sono caratteristiche del \textbf{cammino};
  \item nella \textbf{corsa} non esiste la doppia fase statica: si hanno solo un \textbf{tempo di
    contatto} e un \textbf{tempo di volo}.
\end{itemize}

% ---------------------------------------------------------------------------
\section{Il principio di misura: la legge di Hooke}

Le \textbf{celle di carico} sono costituite da \textit{strain gauge} e seguono la legge di Hooke.

\textbf{Legge di Hooke}: l'allungamento subìto da un corpo elastico è direttamente proporzionale
alla forza a esso applicata; la costante di proporzionalità è detta \textbf{costante elastica} e
dipende dalla natura del materiale.

\rubrica{Dalla deformazione al segnale}
\begin{enumerate}
  \item il gradiente di forza esercitato sullo \textit{strain gauge} ne provoca la
    \textbf{deformazione};
  \item la deformazione provoca una \textbf{variazione di tensione proporzionale alla forza
    applicata};
  \item calibrando la variazione di tensione con la variazione di forza si ottiene una
    \textbf{relazione lineare} $\rightarrow$ ogni variazione sopra la pedana corrisponde a una
    variazione di forza nell'unità di tempo.
\end{enumerate}

\rubrica{Curva forza-allungamento e limiti di misura}
La curva presenta, nell'ordine: \textbf{regione elastica} $\rightarrow$ \textbf{limite di
proporzionalità} $\rightarrow$ \textbf{limite elastico} $\rightarrow$ \textbf{regione plastica}
$\rightarrow$ \textbf{punto di rottura}.
\begin{itemize}
  \item ogni \textit{strain gauge} ha quindi un proprio \textbf{range di misura};
  \item se la forza applicata supera le capacità elastiche dello strumento, questo \textbf{non
    torna allo stato di quiete} e non fornisce più misure affidabili;
  \item esistono perciò \textit{strain gauge} con scale completamente diverse: quelli delle
    piattaforme, per le variazioni di forza esercitate dall'atleta, e quelli impiegati per misurare
    le tensioni degli smottamenti dei terreni nei terremoti.
\end{itemize}

% ---------------------------------------------------------------------------
\section{L'evoluzione dei test di salto}

Le piattaforme di forza nascono tra la fine degli anni Sessanta e l'inizio degli anni Settanta; i
test iniziali prevedevano invece \textbf{sistemi indiretti}.

\begin{enumerate}
  \item \textbf{Sargent} (1921): misura dell'altezza di salto tramite \textit{reach}, cioè la
    differenza tra il punto toccato dal soggetto su una scala millimetrata a muro e il punto più
    alto raggiunto in volo;
  \item \textbf{Abalakov} (1938): con la fettuccia metrica, differenza tra la posizione eretta e la
    posizione raggiunta nel punto più alto del salto;
  \item \textbf{Davies-Rennie} (1968) e \textbf{Cavagna et al.} (1972): impiego della piattaforma
    di forza;
  \item \textbf{Bosco} (1980): utilizzo del \textbf{tempo di volo} con il tappetino a conduttanza;
  \item \textbf{Vertec} (1990): molto più pratico del test di Sargent, di cui condivide il principio.
\end{enumerate}

\rubrica{Limite dei test indiretti}
Non forniscono \textbf{nessuna informazione sulle capacità neuromuscolari} né sul modo in cui esse
si sviluppano durante il salto: si stabilisce fino a che altezza il soggetto è arrivato, ma non
\textbf{come} l'ha sviluppata.

% ---------------------------------------------------------------------------
\section{Il salto sulla piattaforma di forza}

Con la piattaforma si ottengono informazioni su \textbf{tutti i processi neuromuscolari} che
precedono e caratterizzano la prestazione, non solo sul risultato finale.

\subsection{Lo squat jump}

\textbf{Squat jump}: salto eseguito senza alcun contromovimento $\rightarrow$ modalità
\textbf{concentrica di tipo esplosivo}, in cui il muscolo si accorcia senza essersi allungato
preventivamente.

\rubrica{Lettura del tracciato}
\begin{itemize}
  \item lo sviluppo di forza procede fino al momento dello stacco, seguito dal \textbf{tempo di
    volo} e dalla \textbf{fase di ricaduta};
  \item il \textbf{picco di forza in ricaduta è molto più alto} di quello della fase propulsiva
    $\rightarrow$ spiega perché molti infortuni avvengono in atterraggio più che in spinta;
  \item anche la fase propulsiva resta un momento delicato: il sistema neuromuscolare vi è espresso
    sempre al massimo.
\end{itemize}

\rubrica{Dalla forza alle altre grandezze}
\begin{itemize}
  \item la \textbf{velocità} diventa massima quando la forza diventa minima: forza e velocità stanno
    in rapporto quasi inverso, così come velocità e accelerazione, perché $F = m \cdot a$;
  \item il \textbf{picco di velocità} si raggiunge allo stacco; poi la velocità decade linearmente
    per effetto della gravità, fino all'impatto al suolo --- dove il soggetto sperimenta la forza di
    reazione più alta;
  \item l'\textbf{altezza massima} corrisponde al momento in cui la velocità torna a $0$, cioè il
    punto più alto della parabola, che coincide con la \textbf{metà del tempo di volo}.
\end{itemize}

\subsection{I tre tipi di salto a confronto}

Tracciati forza-tempo registrati su piattaforma: a parità di tempo di volo, lo sviluppo della forza
è completamente differente.

\begin{enumerate}
  \item \textbf{concentrico esplosivo} (squat jump): la forza è sviluppata da un accorciamento senza
    allungamento preventivo;
  \item \textbf{con contromovimento} (CMJ): si osserva una \textbf{fase eccentrica} caratterizzata
    da un peso apparente ridotto rispetto al peso corporeo --- la forza di reazione del terreno
    diminuisce perché il soggetto va verso il basso --- seguita da un'impennata fino al
    \textbf{picco di forza all'inizio della fase concentrica}, che poi diminuisce fino al volo;
  \item \textbf{pliometrico}: l'espressione di forza è caratterizzata da un picco che si esprime
    durante la \textbf{sola fase di contatto}, in cui si distinguono ugualmente una fase eccentrica
    e una concentrica.
\end{enumerate}

\rubrica{Naturalezza del gesto}
Il movimento più naturale è quello determinato dal \textbf{ciclo stiramento-accorciamento}; il
movimento puramente concentrico è meno naturale e difficile da realizzare, se non appreso
precedentemente.

\subsection{Le fasi del salto e le grandezze calcolate}

Sequenza delle fasi rilevate dalla piattaforma: $standing \rightarrow eccentrica \rightarrow
concentrica \rightarrow volo \rightarrow atterraggio \rightarrow standing\ still$; per ciascuna si
seguono \textbf{forza verticale}, \textbf{velocità}, \textbf{altezza di salto} e \textbf{potenza}.

\begin{itemize}
  \item nella \textbf{fase eccentrica} la forza registrata è più bassa del peso corporeo;
  \item il \textbf{picco di forza} si osserva quando la velocità è $0$, nel punto più basso, cioè
    nella posizione di accosciata $\rightarrow$ cade all'\textbf{inizio della fase propulsiva}, non
    alla fine;
  \item di conseguenza, nelle prestazioni con ciclo stiramento-accorciamento la capacità di
    \textbf{reclutare unità motorie ed esprimere forza deve avvenire all'inizio del movimento};
  \item raggiunto il picco la forza diminuisce, perché diminuisce l'inerzia dell'oggetto da muovere
    --- il peso del corpo --- mentre la \textbf{velocità cresce}; dal momento in cui la velocità
    raggiunge il proprio picco interviene la forza di gravità;
  \item la velocità così espressa è \textbf{altamente correlata con l'altezza di salto}, raggiunta a
    metà del tempo di volo.
\end{itemize}

\subsection{Applicazione al sollevamento pesi}

Confronto delle curve forza-tempo di \textbf{strappo} e \textbf{girata} a parità di carico.
\begin{itemize}
  \item la parte iniziale di \textbf{forza esplosiva} è sviluppata nella fase di spinta;
  \item il \textbf{tempo di volo} corrisponde al momento in cui il soggetto ``si infila'' sotto il
    bilanciere, per poi spingerlo;
  \item la \textbf{girata} mostra uno sviluppo di forza maggiore rispetto allo strappo.
\end{itemize}

% ---------------------------------------------------------------------------
\section{Dalla forza alla velocità e allo spostamento}

Il passaggio dalla curva forza-tempo alla curva velocità-tempo, e da questa alla curva
spostamento-tempo, avviene per \textbf{integrazione}:
\[ dF = m\,\frac{dv}{dt} \qquad dF\,dt = m\,dv \qquad \int_{t_1}^{t_2} dF\,dt = m\,\Delta v \]

\rubrica{Procedimento di calcolo numerico}
\begin{enumerate}
  \item si sottrae alla forza la \textbf{costante peso}, che rimane tale $\rightarrow$ si ottiene
    l'\textbf{accelerazione};
  \item si calcola l'area sottesa dalla curva con il \textbf{metodo dei parallelepipedi}: l'area di
    ogni rettangolo vale $a \cdot t$;
  \item la \textbf{sommatoria} delle singole aree dà la velocità $\rightarrow$ $v(t)$;
  \item applicando lo stesso procedimento alla curva $v(t)$ si ottiene lo \textbf{spostamento} in
    funzione del tempo;
  \item da $V(t)$ si ricavano infine $h(t)$ e $P(t)$.
\end{enumerate}

\subsection{Altezza di salto dalla velocità di stacco}

Nota la \textbf{velocità allo stacco} $v_0$ (\textit{take off velocity}), misurata con la pedana di
forza, si applica l'uguaglianza tra energia cinetica ed energia potenziale:
\[ E_c = E_p \qquad \tfrac{1}{2}mv^2 = mgh \qquad \Rightarrow \qquad h_{max} = \frac{v_0^2}{2g} \]
La massa si semplifica tra i due membri.

\subsection{Altezza di salto dal tempo di volo}

Procedimento utilizzato dalle \textbf{pedane conduttive}, che misurano il tempo di volo $t_f$:
\begin{itemize}
  \item \textbf{velocità allo stacco}: $V_{to} = \frac{t_f}{2} \cdot g$ --- accelerazione per tempo,
    quindi già un'integrazione;
  \item \textbf{velocità media}: $V_{media} = \frac{t_f}{4} \cdot g$;
  \item \textbf{altezza}: $h = V_{media} \cdot \frac{t_f}{2} = \frac{t_f}{4} \cdot g \cdot
    \frac{t_f}{2}$.
\end{itemize}
\[ h = \frac{t_f^2 \cdot g}{8} \]
Formula di Asmussen e Bonde Petersen, \textit{Acta Physiol Scand}, 1974.

% ---------------------------------------------------------------------------
\section{Le pedane a contatto}

\subsection{Ergojump}

\textbf{Pedana a conduttanza}: non è altro che un \textbf{interruttore}, costituito da barre di
metallo che vanno in contatto tra loro quando il soggetto è sopra la pedana e si staccano quando è
in aria.
\begin{itemize}
  \item collegata a un \textbf{timer}, permette di ottenere \textbf{tempo di volo} e \textbf{tempo
    di contatto};
  \item esiste anche la versione con \textbf{sistema ottico}, a barre, che sfrutta lo stesso
    principio.
\end{itemize}

\subsection{Wi-Jump}

Sistema che riprende le caratteristiche della pedana a conduttanza rendendola \textbf{wireless}:
$MAT \rightarrow microcontrollore \rightarrow Bluetooth \rightarrow smartphone$ o palmare.
\begin{itemize}
  \item consente di disporre di tutta una serie di parametri \textbf{nell'immediatezza};
  \item è stato \textbf{validato} contro pedana di forza e sistema optoelettronico (Annino et al.,
    2016; 2017);
  \item le relazioni ottenute sono lineari e altamente correlate: $R^2 = 0,975$ con $p < 0,0001$
    ($n = 26$) e $R^2 = 0,96$ con $p < 0,0001$ ($n = 19$) $\rightarrow$ la strumentazione risulta
    \textbf{valida dal punto di vista scientifico}.
\end{itemize}

% mazzi: 09
\chapter{L'allenamento della potenza muscolare}

\textbf{Allenamento della potenza muscolare}: prescrizione ricavata dalle curve forza-velocità e
potenza-velocità della dinamometria isoinerziale, dove il riferimento non è più il numero di
ripetizioni ma il \textbf{carico} e la \textbf{potenza espressa}.

% ---------------------------------------------------------------------------
\section{I parametri dell'allenamento neuromuscolare}

Due soli parametri definiscono lo stimolo:
\begin{enumerate}
  \item \textbf{entità del carico}: in percentuale dell'1RM;
  \item \textbf{intensità dello stimolo}: in percentuale della potenza massima ($P_{max}$).
\end{enumerate}

\begin{itemize}
  \item la \textbf{potenza} è la vera intensità dell'allenamento: il numero di ripetizioni smette di
    essere il parametro di riferimento;
  \item una volta determinata la curva forza-velocità, sono noti i valori di potenza di ogni carico
    sotteso alla curva;
  \item il carico si sceglie rispetto al \textbf{picco di potenza} --- a sinistra o a destra ---
    secondo l'esigenza e l'obiettivo.
\end{itemize}

% ---------------------------------------------------------------------------
\section{I parametri di riferimento dei quattro tipi di allenamento}

Ogni tipo di allenamento occupa una porzione diversa delle due curve, individuata da una coppia di
valori: quanto carico (\% 1RM) e quanta potenza (\% $P_{max}$) si deve esprimere con quel carico.

\begin{itemize}
  \item \textbf{forza massima} (\textit{maximal strength training}): carico 70--100\,\% 1RM, potenza
    $>$ 90\,\% $P_{max}$;
  \item \textbf{ipertrofia} (\textit{hypertrophic training}): carico 70--90\,\% 1RM, potenza
    70--80\,\% $P_{max}$;
  \item \textbf{forza esplosiva} (\textit{explosive strength training}): carico 20--70\,\% 1RM,
    potenza $>$ 90\,\% $P_{max}$;
  \item \textbf{resistenza alla forza veloce} (\textit{speed strength endurance training}): carico
    20--50\,\% 1RM, potenza 80--90\,\% $P_{max}$.
\end{itemize}

\rubrica{Perché la potenza cambia il senso del carico}
\begin{itemize}
  \item con carichi $>$ 70\,\% 1RM il reclutamento è quasi massimale: a fine serie tutte le fibre
    risultano esaurite, ma quelle davvero reclutate e portate a esaurimento sono le \textbf{fibre
    veloci};
  \item nell'ipertrofia la potenza si abbassa deliberatamente (fino all'80\,\%) per stressare il
    sistema dal punto di vista metabolico: maggiore produzione di acido lattico $\rightarrow$
    aumento del GH $\rightarrow$ ipertrofia; interessa soprattutto lo stimolo sul
    \textbf{testosterone}, fenotipizzatore delle fibre veloci;
  \item nella forza esplosiva il carico può stare a destra come a sinistra del picco, ma la potenza
    richiesta resta $>$ 90\,\%: l'obiettivo è condizionare \textbf{pattern neuromuscolari}
    specifici, cioè strategie di reclutamento specifiche;
  \item nella resistenza alla forza veloce il carico si sbilancia verso destra (carichi bassi,
    velocità alte) per ottenere uno \textbf{stress metabolico}.
\end{itemize}

Il quadro di riferimento esteso aggiunge la \textbf{resistenza muscolare} e propone per la
resistenza alla forza veloce un intervallo di carico leggermente diverso:

\begin{table}[htbp]
  \centering
  \small
  \renewcommand{\arraystretch}{1.3}
  \begin{tabularx}{\textwidth}{|X|c|c|}
    \hline
    \textbf{Tipo di allenamento} & \textbf{Carico (\% 1RM)} &
      \textbf{\% della $P_{max}$ con quel carico} \\
    \hline
    Forza massima (\textit{maximal strength}) & 70--100\,\% & minimo 90\,\% \\
    \hline
    Forza esplosiva (\textit{speed/strength, power}) & 20--70\,\% & minimo 90\,\% \\
    \hline
    Ipertrofia (\textit{hypertrophy}) & 70--90\,\% & 75--85\,\% \\
    \hline
    Resistenza alla forza veloce (\textit{speed/endurance}) & 30--50\,\% & 80--90\,\% \\
    \hline
    Resistenza muscolare (\textit{muscle endurance}) & 30--70\,\% & 70--85\,\% \\
    \hline
  \end{tabularx}
\end{table}

% ---------------------------------------------------------------------------
\section{Reclutamento delle fibre e recupero tra le serie}

Le tre popolazioni di fibre --- \textbf{ST}, \textbf{FTa}, \textbf{FTb} --- si seguono durante e
dopo lo sforzo distinguendo le fibre \textit{recruited}, \textit{recruited exhausted}, \textit{not
recruited}, \textit{recovered} e \textit{not recovered}.

\rubrica{Durante la serie con carichi $>$ 70\,\% 1RM}
\begin{enumerate}
  \item \textbf{inizio}: sono reclutate le ST e parte delle FTa, le FTb solo in quota minima;
  \item \textbf{durante}: la quota di fibre reclutate ed esaurite cresce e si estende alle FTa e
    alle FTb;
  \item \textbf{fine}: quasi tutte le fibre risultano reclutate ed esaurite.
\end{enumerate}

Lo stesso avviene nella successione di \textbf{sforzi balistici massimali}: anche con carichi bassi
il reclutamento parte dalle fibre veloci e, man mano che si esauriscono, passa alle intermedie e poi
alle lente $\rightarrow$ calo dei valori di potenza.

\rubrica{Il recupero dopo ripetizioni submassimali}
\begin{itemize}
  \item \textbf{subito dopo il lavoro}: la potenza massima esprimibile è nettamente ridotta;
  \item \textbf{dopo 1 minuto}: le ST sono recuperate, ma le FTa risultano ancora \textit{not
    recovered} $\rightarrow$ la potenza non torna al 100\,\%;
  \item \textbf{dopo 3 minuti}: il recupero è più ampio, ma la quota di FTb non recuperate resta il
    fattore che limita la ripresa della serie.
\end{itemize}

Ne segue che il \textbf{tempo di recupero tra le serie} va commisurato alle fibre che interessano:
se sono le veloci, i tempi diventano lunghi e vanno verificati sul campo --- se dopo tre minuti la
potenza esprimibile resta $<$ 90\,\% della massima, il recupero si allunga ancora.

\rubrica{Ipertrofia selettiva}
\textbf{Ipertrofia selettiva}: condizionare solo le fibre che servono alla prestazione, attraverso i
pattern neuromuscolari altamente specifici legati alle \textbf{unità motorie veloci di tipo B}.
\begin{itemize}
  \item si contrappone all'\textbf{ipertrofia generale} del \textit{bodybuilder}, dove l'obiettivo è
    aumentare la massa a prescindere dal pattern neuromuscolare utilizzato;
  \item condizionare unità motorie che non sono il fattore limitante è inutile $\rightarrow$ può
    addirittura disturbare la prestazione finale.
\end{itemize}

I mezzi di allenamento si dispongono lungo un continuum tra \textbf{fenomeni nervosi} e
\textbf{massa muscolare}: i carichi leggeri ad alta velocità e le serie da 1--3RM stanno sul
versante nervoso, il 6RM in posizione intermedia e il 10RM verso la massa muscolare.

% ---------------------------------------------------------------------------
\section{Il monitoraggio della potenza nella serie}

Con l'encoder la potenza media di ogni ripetizione è disponibile in tempo reale. Nell'esempio di una
serie da 10 ripetizioni, la potenza media (W) segue l'andamento:

\begin{center}
246 --- 233 --- 245 --- 229 --- 219 --- 212 --- 196 --- 186 --- 175 --- 188
\end{center}

\begin{itemize}
  \item si fissa una \textbf{soglia di potenza} (per la forza esplosiva $\geq$ 90\,\% della
    $P_{max}$, con carichi indicativamente tra il 30--40\,\% e il 70\,\%);
  \item finché ogni ripetizione resta sopra la soglia si sta lavorando sulle unità motorie veloci;
  \item quando la potenza scende sotto il target è cambiato il pattern di reclutamento, perché
    quello privilegiato si è affaticato $\rightarrow$ la serie si interrompe;
  \item il \textbf{numero ottimale di ripetizioni} non è quindi il massimo eseguibile con quel
    carico, ma il numero di ripetizioni che rientrano nel 90\,\% della potenza massima;
  \item lo stesso criterio individualizza ripetizioni per serie, recupero tra le serie e numero di
    serie: il ``\textit{no pain no gain}'' non vale più, lo sforzo non deve essere massimo ma
    selettivo.
\end{itemize}

Poiché l'adattamento fa crescere le potenze espresse, i valori vanno rimisurati e il carico
rimodulato a ogni progresso: gli allenamenti successivi si tarano sugli adattamenti già ottenuti, e
il volume di lavoro resta quantificato, evitando situazioni di sovraccarico.

% ---------------------------------------------------------------------------
\section{Dal test diagnostico alla programmazione}

Sul principio $P = F \times v$, il \textbf{test diagnostico} (\textit{ergo-power}) alimenta cinque
usi distinti:
\begin{itemize}
  \item \textbf{pianificazione e controllo} dei piani di allenamento per atleti di alta prestazione;
  \item \textbf{pianificazione e controllo} dei carichi di lavoro in riabilitazione;
  \item \textbf{pianificazione} dei carichi per l'attività fisica generale;
  \item \textbf{confronto} con i valori standard della popolazione normale;
  \item \textbf{confronto} con i modelli teorico-matematici e ideali della prestazione.
\end{itemize}

\rubrica{La valutazione della capacità di salto}
Gli stessi parametri dell'encoder (spostamento, velocità media, picco di velocità, tempo, forza
media) si applicano al salto, con in più la \textbf{\textit{jump height}}: la differenza tra la
posizione di partenza e quella di arrivo del soggetto. Per il lavoro con i sovraccarichi restano
però decisivi il carico e la percentuale di potenza esprimibile con quel carico.

% ---------------------------------------------------------------------------
\section{Mezzo squat tradizionale e mezzo squat con discesa veloce}

Il \textbf{mezzo squat} si scompone in quattro momenti, letti sulla velocità:
\begin{enumerate}
  \item \textbf{posizione eretta}: velocità nulla;
  \item \textbf{discesa}: la velocità si abbassa rapidamente (valori negativi);
  \item \textbf{massima accosciata}: la velocità torna a zero;
  \item \textbf{inversione} tra fase eccentrica e concentrica e fase finale.
\end{enumerate}
Tra il punto 3 e il punto 4 si colloca il \textbf{picco di velocità} del sollevamento.

\textbf{Mezzo squat con discesa veloce}: variante in cui si velocizza la fase eccentrica, lasciando
invariato il resto dell'esercizio.
\begin{itemize}
  \item accumula più \textbf{energia elastica} e attiva il \textbf{riflesso da stiramento}
    $\rightarrow$ restituiti nella fase concentrica successiva;
  \item modifica i parametri meccanici dell'esercizio e con essi la risposta neuromuscolare.
\end{itemize}

\rubrica{Che cosa cambia nei tracciati}
Registrando forza verticale, posizione verticale del centro di gravità e velocità nelle due
esecuzioni:
\begin{itemize}
  \item \textbf{velocità}: nell'esecuzione veloce la fase eccentrica raggiunge velocità (negative)
    nettamente maggiori e il picco concentrico successivo è più alto; l'intera prova si compie in un
    tempo più breve;
  \item \textbf{potenza}: il picco è più alto nell'esecuzione veloce e viene raggiunto prima ---
    quello che interessa è la \textbf{pendenza della curva}, cioè la capacità di esprimere potenza
    nel minor tempo possibile;
  \item \textbf{forza}: nell'esecuzione veloce il picco è più alto e anticipato rispetto a quella
    tradizionale, sia nella fase di spinta sia all'impatto nella fase di salto.
\end{itemize}
I due sistemi attivano dunque pattern neuromuscolari completamente diversi.

% ---------------------------------------------------------------------------
\section{L'allenamento a isopotenza}

\textbf{Isopotenza}: lavorare a un livello di potenza costante scegliendo carichi diversi. A uno
stesso livello di potenza, sotto il picco della curva potenza-velocità, corrispondono due carichi
distinti --- uno a sinistra e uno a destra.
\begin{itemize}
  \item permette di mantenere il livello di potenza voluto con il carico più basso, senza
    sovraccaricare la schiena;
  \item è il presupposto metodologico del confronto sperimentale che segue.
\end{itemize}

% ---------------------------------------------------------------------------
\section{Il confronto sperimentale tra carico alto e carico basso}

\rubrica{Metodologia}
\begin{itemize}
  \item \textbf{gruppo di controllo}: 5 serie da 6 ripetizioni con il 70\,\% dell'1RM;
  \item \textbf{gruppo sperimentale}: 5 serie da 6 ripetizioni con un carico pari al 45\,\%
    dell'1RM.
\end{itemize}

\rubrica{Risultati}
Dopo 4 settimane di allenamento entrambi i gruppi migliorano, con differenze tra i gruppi per forza,
potenza e sprint ($p < 0{,}05$).

\begin{table}[htbp]
  \centering
  \small
  \renewcommand{\arraystretch}{1.3}
  \begin{tabularx}{\textwidth}{|l|>{\centering\arraybackslash}X|>{\centering\arraybackslash}X|>{\centering\arraybackslash}X|>{\centering\arraybackslash}X|}
    \hline
    & \textbf{RM (kg)} & \textbf{CMJ (cm)} & \textbf{PPO (W)} & \textbf{20 m (s)} \\
    \hline
    \multicolumn{5}{|c|}{\textbf{TG\textsubscript{FAST}} --- carico 45\,\% 1RM} \\
    \hline
    Pre & 120,8 $\pm$ 12,4 & 31,9 $\pm$ 3,6 & 475,1 $\pm$ 66,4 & 3,61 $\pm$ 0,15 \\
    \hline
    Post & 131,9 $\pm$ 18,0 & 36,0 $\pm$ 4,1 & 530,4 $\pm$ 69,9 & 3,54 $\pm$ 0,12 \\
    \hline
    Diff & 11,1 $\pm$ 14 & 4,1 $\pm$ 1,9 & 55,2 $\pm$ 40,5 & $-0,07 \pm 0,10$ \\
    \hline
    \multicolumn{5}{|c|}{\textbf{TG\textsubscript{TRAD}} --- carico 70\,\% 1RM} \\
    \hline
    Pre & 114,2 $\pm$ 14,5 & 30,5 $\pm$ 2,9 & 428,5 $\pm$ 24,8 & 3,64 $\pm$ 0,24 \\
    \hline
    Post & 120,9 $\pm$ 11,6 & 34,5 $\pm$ 3,2 & 459,5 $\pm$ 36,9 & 3,51 $\pm$ 0,15 \\
    \hline
    Diff & 6,7 $\pm$ 16,4 & 4,0 $\pm$ 1,1 & 31,0 $\pm$ 35,6 & $-0,14 \pm 0,12$ \\
    \hline
  \end{tabularx}
\end{table}

Tutte le differenze pre-post sono statisticamente significative in entrambi i gruppi: allenarsi con
un carico molto più basso, purché eseguito in modo esplosivo, produce miglioramenti dello stesso
ordine di quelli ottenuti con carichi alti. La conseguenza metodologica è che si può modulare carico
e volume in funzione delle esigenze del singolo atleta, privilegiando comunque i movimenti esplosivi
e veloci.

% ---------------------------------------------------------------------------
\section{Individualizzazione e specificità del lavoro}

\textbf{Allenamento individuale}: le metodologie devono ruotare attorno alle esigenze del singolo
giocatore, non il contrario --- nel rispetto della sua salute e delle sue potenzialità.
\begin{itemize}
  \item carico e volume si modulano sulle caratteristiche neuromuscolari e metaboliche del singolo
    atleta;
  \item quando il volume di lavoro va ridotto, si privilegiano gli allenamenti che condizionano in
    modo specifico;
  \item il tempo per il lavoro con i sovraccarichi è poco, a fronte dell'enorme carico
    tecnico-tattico e della partita: va reso il più specifico possibile e spendibile il giorno della
    gara.
\end{itemize}

\textbf{Specificità}: coerenza biomeccanica, neuromuscolare e metabolica con i movimenti che si
ritrovano in partita.
\begin{itemize}
  \item in partita nessuno esegue uno squat con i sovraccarichi: ciò che si condiziona non è il
    gesto, ma il \textbf{pattern neuromuscolare} che serve a sprint, accelerazioni, decelerazioni e
    cambi di direzione ad altissima intensità;
  \item il vantaggio è il \textbf{trasferimento} immediato alla prestazione, senza dover attendere
    ulteriore tempo.
\end{itemize}

\rubrica{Il sovraccarico in prossimità della partita}
Con carichi e volumi adeguati il lavoro con i sovraccarichi può essere collocato anche poco prima
della partita:
\begin{itemize}
  \item non come condizionamento, ma come \textbf{riscaldamento e attivazione} dei pattern
    neuromuscolari che serviranno subito dopo;
  \item è sufficiente una serie di 2--3 ripetizioni: più che appesantire il sistema neuromuscolare,
    lo attiva;
  \item con tutte le attenzioni dovute al \textbf{volume}.
\end{itemize}


\end{document}

% =============================================================================
% PROMEMORIA STILE SCHEMATICO (vedi CONVENZIONI.md, §7)
% =============================================================================
% - Definizioni:      \textbf{Termine}: spiegazione sintetica;   (due punti, non "=")
% - Sotto-titoli:     \rubrica{Struttura}  -> grassetto a riga propria, non numerato
% - Liste annidate:   enumerate 1. -> (a) -> i. ; itemize • -> – ; sono lo strumento base
% - Enfasi:           grassetto = termine chiave; \textit{corsivo} = termine secondario
% - Processo/catena:  $molecole \rightarrow cellule \rightarrow tessuti \rightarrow organi$
% - Diametro:         $\phi$ 7,5\,$\mu$m      (misure e formule in math mode)
% - Formule:          $O_2$, $Na^+$, $H_2O$, $C_6H_{12}O_6$, $10^{-9}s$
% - Albero:           \begin{center}\begin{forest} classalbero [Radice [A][B]] \end{forest}\end{center}
% - Due colonne:      \begin{multicols}{2} ... \end{multicols}   (liste lunghe)
% - Tabelle:          a griglia (\hline + |) e con parsimonia; preferire le liste annidate
